\chapter{Literature Review}
\label{chap:lr}
\section{Fundamentals of Subsurface Fluid Dynamics}
\label{sec:theoretical_foundations}

Reservoir management is a decision-making process under geological uncertainty. Da Cruz and Deutsch~\cite{pangea_uncertainty} note that development decisions, from well placement to production scheduling, rely on subsurface models whose calibration data come almost exclusively from wellbores. Wellbores sample a negligible fraction of the reservoir volume, so the reliability of production forecasts depends on how well the porous medium is characterized between the wells and on the accuracy with which the flow equations are applied.

A hydrocarbon reservoir is a porous geological formation that stores and transmits fluids. Two scalar fields define its state for any hydrodynamic calculation~\cite{ssdynamics_guide}: porosity $\phi$, a dimensionless quantity that ranges from below 0.05 in tight formations to over 0.30 in clean sandstones, and absolute permeability, represented by a second-rank tensor $K$ that describes the conductivity of the connected pore network. Porosity varies gradually in space, whereas permeability is heterogeneous and varies by orders of magnitude over short distances. Such variability reflects deterministic depositional processes rather than stochastic noise. A review of rule-based geological modeling~\cite{subsurfaceai_heterogeneity} describes architectural elements such as high-permeability channels, shale drapes, and inclined heterolithic stratification. These structures are smaller than the typical 10-meter resolution of seismic surveys, and they control reservoir connectivity and sweep efficiency. A predictive model must resolve them to reproduce the non-linear dynamics of fluid displacement.

The structural architecture of a reservoir also controls the flow regime. A standard morphological classification~\cite{studfile_types} distinguishes layered reservoirs, in which properties vary vertically but stay roughly constant laterally and flow follows high-permeability layers, from naturally fractured reservoirs, which form a dual-porosity system: a low-permeability matrix that stores fluids and a high-permeability fracture network that conducts them. In a fractured-porous medium the macroscopic flow follows the connectivity of the fracture network, which produces non-linear flow paths and sharp displacement fronts that coarse-grid discretizations resolve poorly.

Within the classification of fractured-porous formations introduced in Kanevskaya, Kadet, and Selyakov~\cite{textbook_fpor}, Type II reservoirs are those in which porous matrix blocks are fully surrounded by fractures. The blocks do not form a continuous medium, and macroscopic filtration occurs only through the fracture network.

Multiphase flow in such systems is described by a single-level formulation under two assumptions: instantaneous saturation equilibrium within each matrix block, and local pressure equilibrium between a block and the surrounding fractures. The thermodynamic state of the system is then captured by three continuous scalar fields: a global pressure $P$, the water saturation $S$ in the fracture network, and the water saturation $\bar{S}$ in the matrix blocks. Variables with an overbar denote matrix-block quantities throughout this section.

Unlike in conventional porous media, the cross-flow between matrix blocks and fractures is driven mainly by elastic forces and not by capillary pressure. Its intensity $q_\Sigma$ is proportional to the time derivative of the global pressure~\cite{textbook_fpor}:
\begin{equation}
    q_\Sigma = -\bar{\beta} \frac{\partial P}{\partial \tau},
    \label{eq:mass_transfer}
\end{equation}
where $\bar{\beta}$ is the elastic capacity coefficient of the matrix blocks and $\tau$ is continuous time.

Macroscopic transport through the fracture network follows an extension of Darcy's law. In the Type II framework, the total filtration velocity is
\begin{equation}
    \mathbf{V} = - \frac{K K^*}{\mu_1} \nabla P,
    \label{eq:darcy_fpor}
\end{equation}
where $K$ is the absolute permeability tensor of the fracture system, $K^*$ is the effective phase mobility, and $\mu_1$ is the reference fluid viscosity. Combining this convective flow with the elastic drainage from the matrix blocks gives the coupled system of non-linear partial differential equations~\cite{textbook_fpor}:
\begin{align}
    \beta \frac{\partial P}{\partial \tau} + \text{div}\,\mathbf{V} &= q_\Sigma, \label{eq:fpor_pressure} \\
    m \frac{\partial S}{\partial \tau} + \beta_1^* S \frac{\partial P}{\partial \tau} + \text{div}(f \mathbf{V}) &= \lambda \, q_\Sigma, \label{eq:fpor_frac_sat} \\
    \bar{m} \frac{\partial \bar{S}}{\partial \tau} + \bar{\beta}_1^* \bar{S} \frac{\partial P}{\partial \tau} &= -\lambda \, q_\Sigma. \label{eq:fpor_block_sat}
\end{align}
Here $m$ and $\bar{m}$ are the effective porosities of the fractures and the matrix blocks, $f$ is the fractional flow function of water in the macroscopic Darcy flow through the fractures, and $\lambda$ is the fractional composition of the inter-porosity exchange between matrix and fractures. Relative phase permeabilities describe how the mobility of each phase depends on local saturation; in the fractures they depend linearly on saturation, while in the matrix blocks they follow a higher-order power law governed by capillary effects~\cite{textbook_fpor}.

The non-linear system above produces a hydrodynamic instability known as viscous fingering, which is amplified in fractured-porous media: the fracture network provides preferential flow paths. The mobility ratio, defined as the ratio of the mobility of the displacing phase to that of the displaced phase, serves as the stability indicator: when it exceeds unity, small perturbations at the displacement front grow rather than decay~\cite{istina_fingering}. A low-viscosity water phase displacing a high-viscosity oil therefore channels rapidly through the fracture network in irregular finger-like patterns. It bypasses oil trapped in the low-permeability matrix blocks and produces early water breakthrough at the producers. Resolving these fingers numerically requires a fine spatial discretization, since they are narrow features set by the local interplay of permeability and pressure gradients. Coarse grids reduce computational cost but introduce numerical diffusion that smears the sharp gradients and suppresses the instability, biasing recovery forecasts towards optimistic values. Accurate simulation of viscous fingering requires grid resolutions that are usually too fine to be feasible for full-field models.




\section{Numerical Modeling}
\label{sec:numerical_challenges}

The coupled system in Equations~\ref{eq:fpor_pressure}--\ref{eq:fpor_block_sat} has no analytical solution for heterogeneous reservoirs, because of the non-linearity of relative permeability and the complexity of the boundary conditions. The industry therefore relies on numerical approximation, primarily the finite difference method (FDM) and the finite volume method (FVM).

Peaceman~\cite{peaceman1977fundamentals} formalized the spatial discretization: in the block-centered formulation, petrophysical properties and state variables are averaged over the discrete volume of each grid block. Partial derivatives are then approximated by finite difference quotients obtained from Taylor expansions. For a scalar function $u(x)$, the central difference scheme gives
\begin{equation}
    \frac{\partial u}{\partial x} \approx \frac{u(x + \Delta x) - u(x - \Delta x)}{2\Delta x},
    \label{eq:central_diff}
\end{equation}
where $\Delta x$ is the spatial grid spacing.

Discretization introduces a local truncation error $\epsilon_L$. For a second-order central difference scheme~\cite{peaceman1977fundamentals},
\begin{equation}
    \epsilon_L = \mathcal{O}(\Delta x^2) + \mathcal{O}(\Delta t),
\end{equation}
where $\Delta t$ is the time step. Accuracy is tied to grid resolution: halving the grid spacing cuts the approximation error by a factor of four and, in three dimensions, multiplies the number of cells $N$ by eight. Since the cost of solving the resulting linear systems scales super-linearly in $N$, the cost of high-fidelity simulation becomes impractical at full-field scale.

Generalizations of the Type II model add to the computational load. While Type II isolates macroscopic flow to the fracture network, generalized dual-permeability models~\cite{konyukhov2019numerical, volkov2016numerical} describe simultaneous filtration in both the matrix and the fractures, which is the Type I configuration. The discrete pressure equation for these coupled systems is parabolic and is solved by implicit difference schemes. In the notation of Konyukhov et al.~\cite{konyukhov2019numerical} it takes the form
\begin{equation}
    \mu_3 (\alpha_T + \bar{\alpha}_T) \frac{\partial P}{\partial t} \approx \text{div}\,(\nabla P + \rho g),
    \label{eq:pressure_discrete}
\end{equation}
where $\alpha_T$ and $\bar{\alpha}_T$ are the transport coefficients of the fracture and matrix systems that incorporate both absolute permeability and the effective phase mobility of each medium, $\mu_3$ is the reference viscosity, $\rho$ is the fluid density, and $g$ is gravitational acceleration. The presence of two transport coefficients on the left-hand side captures the additional coupling that must be resolved when both media carry macroscopic flow.

The underlying trade-off between resolution and cost is unchanged by algorithmic improvements. Resolving fine-scale instabilities such as viscous fingering requires grid spacing finer than the scale of heterogeneity and leads to models with billions of cells. Coarsening the grid raises the truncation error and introduces artificial numerical diffusion. In fractured-porous media this diffusion smears the sharp saturation fronts in the fractures and dampens the local pressure gradients. Since the inter-porosity exchange in Equation~\ref{eq:mass_transfer} is driven by the time derivative of pressure, dampened gradients corrupt the mass-transfer calculation and degrade the recovery forecast. This is the trade-off addressed by the learning-based surrogate models examined in Section~\ref{sec:dl_architectures}.


\section{Performance Metrics for Flow Reconstruction}
\label{sec:metrics}

The quantitative assessment of surrogate models for multiphase reservoir flow requires more than a single point-wise score. Non-linear displacement fronts have sharp spatial profiles, and a metric that averages errors over the grid can hide topological errors that change the engineering interpretation of the result. The literature on super-resolution and on remote sensing provides a layered set of metrics that target, in turn, point-wise accuracy, structural geometry, radiometric consistency across channels, the angular alignment of multi-channel signals, and physics-informed properties of the reconstructed field.

The Mean Absolute Error and the Mean Squared Error are the two basic point-wise measures. For a reference field $X$ and a predicted field $Y$ defined on $N$ grid cells,
\begin{equation}
    \text{MAE}(X, Y) = \frac{1}{N} \sum_{i=1}^{N} |x_i - y_i|, \qquad
    \text{MSE}(X, Y) = \frac{1}{N} \sum_{i=1}^{N} (x_i - y_i)^2.
\end{equation}
Mathieu et al.~\cite{mathieu2016deep} argued that MSE is biased towards spatially averaged solutions, and that MAE, being less sensitive to large outliers, tends to preserve edges better in dense prediction tasks. The Peak Signal-to-Noise Ratio rewrites MSE on a logarithmic scale relative to the dynamic range of the reference signal~\cite{gonzalez2018digital}:
\begin{equation}
    \text{PSNR}(X, Y) = 10 \cdot \log_{10}\!\left(\frac{R(X)^{\,2}}{\text{MSE}(X, Y)}\right),
\end{equation}
where $R(X) = \max(X) - \min(X)$ is the dynamic range of the reference field $X$. The dynamic range adapts the metric to the actual span of values in each channel, which is needed when channels follow different normalization schemes or have different numerical ranges. PSNR is reported per channel and as a channel average. A complementary indicator of robustness is the maximum absolute error, $\max_i |x_i - y_i|$, which characterizes the worst-case point deviation and is sensitive to localized non-physical spikes.

Point-wise metrics describe per-cell deviations but do not capture the geometry of fluid fronts or the connectivity of swept regions. Wang et al.~\cite{wang2004image} introduced the Structural Similarity Index (SSIM) to separate the similarity assessment from per-pixel intensity comparison. SSIM decomposes the comparison into three statistical components computed over a local sliding window: a luminance term $l$, a contrast term $c$, and a structural cross-correlation term $s$,
\begin{align}
    l(x, y) &= \frac{2\mu_x\mu_y + C_1}{\mu_x^2 + \mu_y^2 + C_1}, &
    c(x, y) &= \frac{2\sigma_x\sigma_y + C_2}{\sigma_x^2 + \sigma_y^2 + C_2}, &
    s(x, y) &= \frac{\sigma_{xy} + C_3}{\sigma_x\sigma_y + C_3},
\end{align}
with local means $\mu_x, \mu_y$, standard deviations $\sigma_x, \sigma_y$, cross-covariance $\sigma_{xy}$, and small numerical constants $C_1, C_2, C_3$ that stabilize the division near zero. The composite index used in practice combines luminance and the joint contrast-structure term,
\begin{equation}
    \text{SSIM}(x, y) = \frac{(2\mu_x\mu_y + C_1)(2\sigma_{xy} + C_2)}{(\mu_x^2 + \mu_y^2 + C_1)(\sigma_x^2 + \sigma_y^2 + C_2)},
\end{equation}
and takes values in $[-1, 1]$, with $1$ corresponding to perfect structural identity. SSIM penalizes the artificial smoothing of sharp fronts and the geometric misalignment of flow barriers more heavily than uniform offsets, which makes it informative for fields that contain narrow displacement structures.

Multi-channel physical states require evaluation criteria that respect the joint structure of the channels. Wald~\cite{wald2002data} introduced the Error Relative Global Adimensional Synthesis (ERGAS) in the context of multispectral image fusion. For a reference field $X$ and a predicted field $Y$ with $C$ channels,
\begin{equation}
    \text{ERGAS}(X, Y) = 100 \, \frac{h}{l} \, \sqrt{\frac{1}{C} \sum_{c=1}^{C} \frac{\text{MSE}_c}{\mu_c^2}},
\end{equation}
where $h/l$ is the spatial resolution ratio between the high- and low-resolution domains, $\text{MSE}_c$ is the mean squared error of the $c$-th channel, and $\mu_c$ is the mean of the $c$-th reference channel. The per-channel normalization by $\mu_c^2$ makes ERGAS appropriate when channels span different numerical ranges, as for coupled pressure and saturation fields. Kruse et al.~\cite{kruse1993spectral} introduced a complementary criterion, the Spectral Angle Mapper (SAM), which evaluates the preservation of inter-channel relationships by computing the angle between the reference and predicted channel vectors at each spatial location,
\begin{equation}
    \text{SAM}(x, y) = \arccos\!\left( \frac{\sum_{c=1}^{C} x_c y_c}{\sqrt{\sum_{c=1}^{C} x_c^2} \, \sqrt{\sum_{c=1}^{C} y_c^2}} \right).
\end{equation}
A low SAM value indicates that the surrogate model preserves the joint statistics across channels rather than reconstructing them independently. For a dual-porosity formulation, in which pressure, fracture saturation, and matrix saturation are coupled through the inter-porosity exchange term of Section~\ref{sec:theoretical_foundations}, a low SAM corresponds to a reconstruction that respects the coupling and not merely the marginal distribution of each channel.

A separate family of metrics, originating in image restoration, targets the geometry of edges and fronts rather than the global statistics of the field. Spatial gradients computed with the orthogonal Sobel operators~\cite{sobel1968isotropic} approximate these local edge structures, and Mathieu et al.~\cite{mathieu2016deep} formulated image-prediction objectives based directly on such gradients. A gradient-consistency criterion computes the magnitude of the Sobel gradient on both the reference and the reconstructed field and reports the mean absolute difference of the two magnitudes. Unlike MSE, this criterion penalizes the smearing of fronts but tolerates small spatial misalignments, which makes it sensitive to the recovery of front sharpness rather than to absolute front position.

A further family of criteria embeds physical priors directly into the evaluation. Raissi et al.~\cite{arxiv_raissi_pinns} formalized this idea in the physics-informed neural network framework, where residuals of the governing partial differential equation are used as a training signal that constrains the network output to obey the underlying physics. The same principle applies to evaluation: for diffusive fields governed by parabolic equations, such as the pressure field in the system of Equations~\ref{eq:fpor_pressure}--\ref{eq:fpor_block_sat}, the magnitude of a discrete Laplacian operator on the reconstructed field measures the deviation from the smoothness implied by the governing equation. Combined with the gradient-based criterion above, such physics-aware indicators provide a domain-specific layer of evaluation that complements purely statistical metrics by reporting to what extent a reconstruction departs from properties enforced by the governing equations.


\section{Evolution of Super-Resolution Models}
\label{sec:dl_architectures}

The mathematical reconstruction of high-fidelity physical state fields from coarse-grid numerical approximations requires predictive models capable of capturing complex spatial hierarchies and non-linear physical dependencies. While multiphase reservoir simulation is a transient and time-dependent process, the continuous temporal evolution of the system can be discretized into a sequence of spatial snapshots. Consequently, within the framework of machine learning, the acceleration of the overall simulation can be formulated as a spatial super-resolution task applied sequentially to each discrete computational time step. To adapt classical computer vision architectures to the transient physics of fractured-porous media, a specific multi-variable tensor representation is required. Instead of standard multi-color images, the thermodynamic state of the reservoir at any given temporal snapshot is mapped into a physical three-channel tensor denoted by $X \in \mathbb{R}^{3 \times H \times W}$. These three channels correspond precisely to the fundamental state variables of the Type II dual-porosity model: global pressure, fracture saturation, and matrix block saturation. Convolutional Neural Networks have established themselves as the standard mathematical framework for processing such highly structured spatial data. Their efficacy stems from the ability to extract hierarchical features through the sequential application of learnable convolution kernels, allowing the network to abstract localized spatial variations into global semantic representations.

\subsection{Convolutional and Residual Networks}
\label{subsec:dl_foundations}

Prior to the advent of deep learning, spatial upscaling in computational fluid dynamics relied almost exclusively on deterministic interpolation methods, such as bilinear or bicubic filtering. While computationally inexpensive, these deterministic low-pass filters fail to inject new high-frequency information, resulting in the artificial smearing of sharp displacement fronts and exacerbating numerical diffusion. 

The application of deep learning to the spatial super-resolution domain was pioneered by Dong et al.~\cite{dong2014learning} with the introduction of the Super-Resolution Convolutional Neural Network. This architecture reinterpreted the classical sparse-coding pipeline as a deep convolutional network, where the non-linear mapping between low-resolution and high-resolution patch spaces is learned end-to-end. The network consists of three distinct operations: patch extraction and representation using $9 \times 9$ kernels, non-linear feature mapping via $1 \times 1$ or $5 \times 5$ convolutions, and final spatial reconstruction. By utilizing a rectified linear activation function, the architecture achieved superior reconstruction accuracy compared to traditional interpolation. However, the model necessitates the preliminary upsampling of the input low-resolution tensor using bicubic interpolation before any feature extraction occurs. This pre-upsampling strategy forces the network to perform all convolution operations in the high-dimensional space. As a result, the computational complexity scales quadratically with the upscaling factor, causing substantial computational redundancy and limiting the practical depth of the network to only three layers. Such a shallow receptive field heavily restricts its ability to capture the long-range spatial dependencies characteristic of complex geological fracture networks.

To resolve the computational bottlenecks inherent to pre-upsampling, Shi et al.~\cite{shi2016real} introduced the Efficient Sub-Pixel Convolutional Neural Network. The fundamental innovation of this model was the proposition that all feature extraction layers should operate exclusively in the low-resolution space. To upscale the final feature maps back to the high-resolution domain, the authors introduced an efficient sub-pixel convolution layer, frequently referred to as the PixelShuffle operator. Instead of relying on predefined interpolation, this layer learns an array of upscaling filters and employs a periodic shuffling operator to rearrange the elements of a multi-channel low-resolution tensor into a single-channel high-resolution output. By shifting the computational burden to the low-dimensional space, the sub-pixel architecture significantly reduced both memory consumption and inference time, establishing post-upsampling as the definitive standard for modern super-resolution models.

While sub-pixel convolution solved the computational efficiency problem, the quest for higher reconstruction accuracy required deeper network architectures to expand the receptive field. However, simply stacking more layers resulted in the degradation problem, where accuracy saturates and then degrades rapidly. He et al.~\cite{he2016deep} addressed this fundamental obstacle with the proposal of Residual Networks. They hypothesized that it is easier for a network to optimize a residual mapping, formulated as $F(x) := H(x) - x$, rather than the original unreferenced mapping $H(x)$. To realize this, the architecture introduces shortcut connections that perform identity addition, bypassing one or more non-linear convolution layers. This additive formulation allows gradients to propagate directly through the network during backpropagation, effectively solving the vanishing gradient problem and enabling the training of very deep networks. Importantly, when adapting the residual architecture for super-resolution tasks, downsampling operations such as average pooling or max pooling are removed. By maintaining a constant spatial resolution throughout the residual blocks, the network avoids the substantial loss of high-frequency spatial details required for accurate field reconstruction.

Simultaneously, Kim et al.~\cite{kim2016accurate} proposed the Very Deep Super-Resolution network, significantly increasing the network depth to twenty convolutional layers. Unlike the aforementioned residual architecture, which utilizes local shortcut connections block-by-block, this approach introduced the concept of global residual learning. The internal architecture consists of a straightforward chain of convolutions to extract deep features. However, rather than predicting the final state directly, the network explicitly learns only the global residual field, which represents the high-frequency difference between the interpolated low-resolution input and the high-resolution ground truth. This architectural expansion extended the receptive field from $13 \times 13$ to $41 \times 41$ grid cells, enabling the model to exploit contextual information from larger spatial regions. Furthermore, Kim et al. introduced adjustable gradient clipping, which allowed the model to achieve rapid convergence using learning rates four orders of magnitude higher than previous shallow networks. To address the spatial reduction issue prevalent in earlier unpadded networks, the authors systematically applied zero-padding before every convolutional layer to preserve the exact spatial dimensions of the feature maps throughout the forward pass. In the context of reservoir modeling, mitigating boundary shrinkage is necessary to preserve physical boundary conditions at the edges of the simulation grid. However, despite these advancements, this very deep network retained the computationally expensive pre-upsampling strategy.

To address the limitations of pre-upsampling while maximizing the benefits of network depth, Lim et al.~\cite{lim2017enhanced} introduced the Enhanced Deep Super-Resolution and Multi-scale Deep Super-Resolution architectures. These models effectively merged the depth of residual networks with the computational efficiency of the sub-pixel post-upsampling operator. A key architectural modification in these models is the complete removal of Batch Normalization layers from the standard residual block. While batch normalization stabilizes training in classification tasks by standardizing mini-batch statistics, Lim et al. demonstrated that it disrupts the absolute range of continuous state variables in regression problems. For physical multiphase systems, where global pressure and local saturations operate on vastly different numerical scales, batch-wise normalization artificially distorts cross-channel physical correlations. By eliminating this normalization, the enhanced architecture maintains absolute magnitude fidelity while saving approximately forty percent of GPU memory during training. To ensure training stability in such a deep architecture without batch normalization, residual scaling with a factor of 0.1 was adopted. Furthermore, the Multi-scale Deep Super-Resolution architecture incorporates scale-specific processing modules at the input and output while sharing approximately 98 percent of its parameters in a common residual body. This design enables the model to reconstruct target fields of various upscaling factors, such as two, three, and four, simultaneously within a unified framework. As highlighted in recent benchmarks on non-linear diffusion problems~\cite{konyukhov_article}, this multi-scale, unnormalized approach establishes a robust and physically consistent baseline for recovering detailed fields in heterogeneous media.

The lightweight counterpart of the EDSR family is provided by the Modified SRResNet (MSRResNet) architecture, originally introduced by Ledig et al.~\cite{cvf_srgan} as the generator backbone of the SRGAN framework and subsequently refined as a standalone regression model. MSRResNet preserves the core idea of stacked residual blocks operating in low-resolution feature space, followed by sub-pixel upsampling at the network tail, but reduces the number of residual blocks compared to the deeper EDSR configuration. This design provides a favourable trade-off between reconstruction quality and inference cost, which is particularly relevant for interactive surrogate applications in which a single forward pass must complete within a small fraction of a second. As a consequence, MSRResNet is widely adopted as a strong regression baseline in subsequent generative super-resolution research, including the ESRGAN family discussed in Section~\ref{subsec:gan_revolution}.

A parallel research direction focuses on lightweight super-resolution architectures suitable for inference on resource-constrained hardware. Liu et al.~\cite{liu2020rfdn} introduced the Residual Feature Distillation Network (RFDN), which refines the information distillation paradigm originally proposed by Hui et al.~\cite{hui2019imdn}. The RFDN architecture relies on residual feature distillation blocks, each composed of multiple shallow feature distillation paths combined by a $1 \times 1$ convolution and followed by an enhanced spatial attention module. Information distillation isolates the most informative features at each depth through repeated $1 \times 1$ projections, thereby achieving a reconstruction quality comparable to that of much deeper residual networks at a fraction of the parameter count. In the context of reservoir surrogate modeling, this property is essential, since the surrogate must support batch inference over hundreds of temporal snapshots within an interactive desktop environment.

Despite their architectural efficiency, enhanced residual networks are optimized using point-wise loss functions, primarily Mean Squared Error. By minimizing the average magnitude difference between the predicted tensor and the ground truth, these models suffer from a regression-to-the-mean effect. Consequently, while the multi-scale residual architecture provides a robust structural baseline, the underlying problem of high-frequency texture loss remains unresolved. Point-wise optimization guarantees that sharp fluid displacement fronts and highly conductive fracture boundaries will be artificially smoothed, confirming the necessity for a radically different optimization paradigm.

\subsection{Generative Adversarial Frameworks}
\label{subsec:gan_revolution}

The reliance on point-wise optimization metrics, specifically the Mean Squared Error, limits the capacity of deep residual architectures to reconstruct high-frequency spatial gradients. A standard result in statistical estimation theory is that optimizing a neural network parameterized by $\theta$ to minimize the expected squared difference between the predicted tensor $G_\theta(X)$ and the ground truth $Y$ converges to the conditional expectation of the target distribution, defined as:
\begin{equation}
    G_\theta(X) = \mathbb{E}[Y | X].
\end{equation}
For highly non-linear spatial processes exhibiting multiple valid high-frequency states for a single low-resolution input, this conditional expectation acts as a deterministic spatial average. Consequently, the regression-to-the-mean effect dominates the reconstruction. The network recovers the low-frequency topological structures but fails to synthesize the sharp spatial boundaries characteristic of complex phenomena, necessitating a radical shift in the optimization paradigm.

To address the limitations of deterministic spatial averaging, Ledig et al.~\cite{cvf_srgan} introduced the Super-Resolution Generative Adversarial Network framework. This approach reformulates the reconstruction task from a pure regression problem into a minimax two-player game involving a generative network $G$ and a discriminative network $D$. The generator $G$, parameterized by $\theta_G$, learns to map a low-resolution input tensor $I^{LR}$ to a high-resolution space. Simultaneously, the discriminator $D$, parameterized by $\theta_D$, acts as a dynamic loss function trained to distinguish between the generated tensor $G(I^{LR})$ and the authentic high-resolution ground truth $I^{HR}$. The optimization objective is defined by the following adversarial value function $V(D, G)$:
\begin{equation}
    \min_{\theta_G} \max_{\theta_D} V(D, G) = \mathbb{E}_{I^{HR} \sim p_{train}(I^{HR})}[\log D_{\theta_D}(I^{HR})] + \mathbb{E}_{I^{LR} \sim p_{G}(I^{LR})}[\log(1 - D_{\theta_D}(G_{\theta_G}(I^{LR})))].
\end{equation}
Through this adversarial training procedure, the discriminator penalizes artificial smoothing and pushes the probability distribution of the generated spatial fields to converge toward the true high-frequency distribution of the reference data.

A central component of this adversarial framework is the complete abandonment of point-wise tensor comparison in favor of structural evaluation within a deep hidden representation. Ledig et al. proposed the Perceptual Loss, which utilizes a pre-trained deep convolutional network, specifically the VGG architecture, as an invariant feature extractor. Instead of measuring differences in the raw data space, the metric computes the Euclidean distance between the activation maps extracted from the intermediate layers of the VGG network. For a given spatial layer $j$ producing a feature map of dimensions $W_j \times H_j$, the perceptual loss $\mathcal{L}_{VGG}$ is formulated as:
\begin{equation}
    \mathcal{L}_{VGG} = \frac{1}{W_j H_j} \sum_{x=1}^{W_j} \sum_{y=1}^{H_j} \left( \phi_j(I^{HR})_{x,y} - \phi_j(G_{\theta_G}(I^{LR}))_{x,y} \right)^2.
\end{equation}
Here, the non-linear operator mapping the spatial fields into the deep feature space is denoted by $\phi_j$. This formulation favours the recovery of complex structural patterns and continuous high-frequency boundaries, successfully restoring the topological fidelity of the generated tensor even if the traditional point-wise accuracy metrics degrade.

Building upon these foundational adversarial concepts, Wang et al.~\cite{wang2018esrgan} developed the Enhanced Super-Resolution Generative Adversarial Network, which currently serves as a state-of-the-art baseline for rigorous scientific applications. This architecture introduced two profound structural modifications to stabilize training and improve the recovery of continuous state variables. First, the standard residual units were replaced with Residual-in-Residual Dense Blocks to maximize the network capacity without excessively increasing the parameter count. Second, incorporating the insights established by the EDSR architecture discussed in Section~\ref{subsec:dl_foundations}, Wang et al. systematically eliminated all Batch Normalization layers from the generator. As previously demonstrated, avoiding batch-wise standardization is necessary to prevent the distortion of absolute thermodynamic magnitudes across different state channels. Furthermore, Wang et al. enhanced the adversarial component by implementing a Relativistic average Discriminator. Rather than predicting the absolute probability that a given field is real, the relativistic discriminator $D_{Ra}$ estimates the probability that the ground truth tensor is relatively more realistic than the generated counterpart. The relativistic objective is expressed as:
\begin{equation}
    D_{Ra}(I^{HR}, I^{SR}) = \sigma(C(I^{HR}) - \mathbb{E}_{I^{SR}}[C(I^{SR})]),
\end{equation}
where $C$ denotes the non-transformed discriminator output and $\sigma$ represents the standard sigmoid activation function. This relative formulation produces sharper spatial gradients and prevents mode collapse during the optimization process.

The adversarial framework substantially improves structural sharpness, but it also introduces high-frequency artefacts that are not penalized by the discriminator alone. Recovering this control requires complementary signals, either in the form of physics-aware regularization terms acting on the generator output, as discussed in Section~\ref{sec:physics_sr}, or in the form of auxiliary conditioning that constrains the space of plausible reconstructions, as discussed in the next subsection.

\subsection{Conditioning Mechanisms for Multi-task Generators}
\label{subsec:conditioning}

While the architectural innovations described above improve reconstruction of a fixed input modality, the application of super-resolution to physical simulation introduces an additional dimension of variability: identical coarse-grid pressure and saturation fields may correspond to several distinct fine-grid realizations, depending on the operational regime of the well and the instantaneous state of the surrounding reservoir. To resolve this ambiguity, the generator must be conditioned on auxiliary scalar information that complements the spatial input.

A general-purpose conditioning mechanism for deep convolutional networks was introduced by Perez et al.~\cite{perez2018film}, who formalized Feature-wise Linear Modulation (FiLM). FiLM applies an affine transformation to the activations of an intermediate convolutional layer, where the scale and shift parameters are produced by an auxiliary network from the conditioning vector. Formally, for a feature map $F_c$ of channel $c$ and conditioning vector $\mathbf{z}$, the modulated activation is
\begin{equation}
    \widetilde{F}_c = \gamma_c(\mathbf{z}) F_c + \beta_c(\mathbf{z}),
\end{equation}
where the functions $\gamma_c(\cdot)$ and $\beta_c(\cdot)$ are jointly learned with the main network. The same modulation is applied uniformly across spatial positions of each channel, which keeps the parameter count low while injecting global information into every block of the backbone. Perez et al. showed that FiLM provides a robust conditioning mechanism for visual reasoning, and the technique has since been adopted in multi-task super-resolution and image-to-image translation models.

In the context of fractured-porous reservoir simulation, FiLM provides a natural mechanism for supplying the super-resolution generator with the current dynamic state of the simulation, including injection pressure, cumulative water cut, and well operational regime. As discussed in the context of the dual-porosity governing system in Section~\ref{sec:theoretical_foundations}, these quantities directly govern the inter-porosity exchange term and the local mobility ratio, and yet they cannot be reliably inferred from a single coarse-grid snapshot.

The scalar information available alongside a Type II simulation snapshot falls into three structurally different groups. The first group consists of dynamic quantities that evolve at every time step and characterize the instantaneous operating state of the reservoir: the model time, the cumulative water cut of the system as a whole and separately of the fracture and matrix sub-systems, the bottom-hole pressure, the current fluid and oil flow rates of the well, the running oil recovery factors of the fracture and matrix networks, and the residuals of the discrete mass-balance relations used by the numerical solver. The second group is static within a single simulation and describes the physical setting of the reservoir as a whole: the densities, viscosities, heat capacities and thermo-baric coefficients of oil, water, and dissolved gas, the geometry of the reservoir and the well, the operating regime of the well at the outer boundary, the temporal configuration of the run, the size of the computational grid, the regime flags and the tolerances of the numerical solver. The third group describes the layered structure of the medium: each of the layers carries its own geometric thickness and vertical discretization, the porosities of the matrix and the fracture sub-systems, the irreducible and maximum water saturations in both sub-systems, the absolute permeabilities of matrix and fracture, and the activation flags that determine whether the layer participates in inflow through the well or in the boundary inflow. The three groups occupy distinct positions in the governing system of Section~\ref{sec:theoretical_foundations}: the layer-wise quantities parameterise the constitutive coefficients of Equations~\ref{eq:fpor_pressure}--\ref{eq:fpor_block_sat}, the static quantities fix the surrounding thermodynamic and operational regime, and the dynamic quantities track the current state of the system along the trajectory induced by these coefficients. None of the three groups is fully recoverable from a single coarse-grid snapshot, which is what makes them natural targets for an external conditioning mechanism such as FiLM.

The resulting conditioned generator is the basis of the architecture proposed in this thesis and is evaluated against an unconditional baseline in the experimental study.

The adversarial framework, augmented with perceptual evaluation and feature-wise conditioning, therefore provides a coherent answer to two of the limitations identified in Section~\ref{subsec:dl_foundations}: it recovers the high-frequency content lost to point-wise optimization, and it admits external scalar information that disambiguates plausible reconstructions. Whether this framework remains the most appropriate choice in the presence of architectures based on attention, recurrence, and diffusion is examined in the next subsection.

\subsection{Spatiotemporal and Attention-Based Alternatives}
\label{subsec:modern_frontier}

Beyond convolutional and generative-adversarial architectures, three alternative families have been proposed for spatial reconstruction tasks: spatiotemporal recurrent networks, attention-based transformer models, and denoising diffusion models. Each family resolves a specific limitation of the convolutional generator, and a brief comparative assessment is required before justifying the choice of the adversarial framework as the baseline architecture in this work.

Convolutional Long Short-Term Memory networks were introduced by Shi et al.~\cite{shi2015convolutional} to model spatiotemporal correlations through convolutional state transitions. In principle, such an architecture matches the temporal structure of the dual-porosity formulation, since the inter-porosity exchange term in Equation~\ref{eq:mass_transfer} depends on the temporal derivative of pressure. However, autoregressive inference over sequences of hundreds of simulation steps accumulates approximation error at each step. For long-horizon reservoir forecasts, this compounded drift eventually leads to numerical divergence, which contradicts the requirement of stable behaviour across the full simulation history.

Vision Transformer architectures, exemplified by SwinIR~\cite{liang2021swinir}, replace localized convolution by self-attention over windows of spatial tokens. The infinite effective receptive field provided by attention is theoretically well suited to capturing long-range pressure propagation through a highly conductive fracture network. In practice, however, the dense projection operations required by the transformer embedding produce a parameter count and a floating-point cost incompatible with the interactive inference budget targeted in this thesis.

Denoising Diffusion Probabilistic Models~\cite{saharia2021image, rombach2022high} generate high-frequency content by reversing a learned noise process over many iterations. Diffusion models provide the strongest theoretical guarantees for distributional coverage and would be the natural candidate for ensemble-style uncertainty quantification. Their drawback is the iterative sampling procedure, which requires repeated evaluation of the underlying network for each spatial reconstruction. Multiplied by the number of temporal snapshots in a simulation, this cost negates the acceleration that a surrogate model is expected to provide.

These three families occupy complementary niches: recurrent networks optimize temporal memory, transformers resolve global spatial topology, and diffusion models cover the conditional posterior. None of them simultaneously satisfies the constraints of a single-pass, deterministic, interactive surrogate operating on consumer hardware. Within these constraints, the conditioned adversarial framework discussed in Sections~\ref{subsec:gan_revolution} and~\ref{subsec:dl_foundations} provides the most suitable architectural baseline and is adopted as the foundation of the methodology developed in the next chapter.


\section{Super-Resolution in Reservoir Engineering}
\label{sec:physics_sr}

The application of super-resolution techniques to subsurface hydrodynamics necessitates a fundamental redefinition of the reconstruction task. While standard computer vision assumes that low-resolution data is a deterministic and downsampled version of a high-resolution original, scientific super-resolution for partial differential equations operates under a different paradigm. In computational physics, the low-resolution input represents an authentic realization of the governing equations evaluated on a coarse mesh. Such solutions contain numerical diffusion and lack subgrid-scale phenomena, meaning the mapping learned by the neural network involves recovering lost physical information rather than simply inverting a blurring kernel.

Adversarial optimization has proven effective at recovering the fine-scale structures essential for hydrodynamic accuracy. An important validation of this approach was provided by Deng et al.~\cite{deng2019superresolution}, who demonstrated that generative adversarial networks successfully restore the high-frequency energy density of velocity fields. To evaluate the unsteady characteristics of dominant scale eddies, the researchers utilized a two-point spatial correlation coefficient of the velocity components, defined as:
\begin{equation}
    R_{vv}(x, y, x_0, y_0) = \frac{\langle v(x, y) v(x_0, y_0) \rangle}{v_{rms}(x, y) v_{rms}(x_0, y_0)}.
\end{equation}
In this expression, the angle brackets denote the ensemble average over multiple samplings. Their results confirmed that while conventional regression models act as low-pass filters that suppress small-scale fluctuations, the adversarial framework restores high-wavenumber content. This spectral recovery allows for the reconstruction of sharp boundaries and instabilities, such as viscous fingering patterns, which are typically lost to the averaging effects of point-wise optimization.

To reconcile the generative capabilities of adversarial training with the requirements of fluid mechanics, several studies have proposed hybrid architectures that incorporate physical constraints as auxiliary regularization terms in the loss function. A representative example was provided by Doloi et al.~\cite{arxiv_doloi_srcsgan}, who developed a super-resolution constrained generative adversarial network for upscaling oil saturation maps and reported that the standard adversarial objective alone does not guarantee adherence to fluid conservation. Their work demonstrated that augmenting the generator loss with a physically motivated regularization term substantially improves the consistency of the reconstructed saturation field. The principle of complementing data-driven adversarial training with explicit physical regularization, established in their study for single-porosity media, is the same principle that underlies the methodology developed in the present thesis. In contrast to their single-criterion formulation, the present work targets fractured-porous media of Type II and therefore introduces two coupled regularization terms, one acting on the diffusive structure of the pressure channel and the other acting on the geometric coherence of the saturation fronts, as defined in Section~\ref{sec:metrics}.

A complementary line of research focuses on multi-objective optimization strategies for high-contrast channelized media. Aslam, Li, and Yan~\cite{li2025spe} proposed a hybrid multiscale framework that combines the standard pixel-wise error with auxiliary terms penalizing prediction error at well locations, deviation of spatial gradients, total spatial variation, and the local diffusion residual. By processing dynamic flow information alongside static permeability maps, this composite objective achieved average root mean square errors of approximately 76 psi for pressure fields and 0.046 for saturation distributions, while delivering a sixteen-fold computational speed-up relative to fine-scale numerical solvers. These results confirm that combining standard regression with carefully chosen physics-aware regularizers can simultaneously correct coarse-grid inaccuracies and preserve sharp flow boundaries.

Despite the steady progress in deep super-resolution for single-porosity reservoirs, a clear gap remains regarding naturally fractured formations of Type II. To the best of our knowledge, no published study has addressed the simultaneous spatial reconstruction of the three coupled state variables required by the dual-porosity formulation, namely global pressure, fracture saturation, and matrix block saturation. Existing physics-informed frameworks, including those of Doloi et al.~\cite{arxiv_doloi_srcsgan} and Aslam et al.~\cite{li2025spe}, target single-continuum models and rely on conservation criteria adapted to that setting. The methodology developed in the present thesis extends the physics-informed adversarial paradigm to a three-channel tensor representation of the dual-porosity state and introduces two complementary physics-informed regularization terms that act on the pressure field and on the saturation channels respectively. The conditioning of the generator on dynamic simulation scalars further allows the network to disambiguate reconstructions that share an identical coarse-grid input but differ in their underlying operational regime.


\section{Conclusions of the Literature Review}
\label{sec:lr_summary}

The literature surveyed in this chapter motivates the methodological choices adopted in the remainder of this thesis along four dimensions.

First, the dual-porosity formulation of Type II reservoirs, presented in Section~\ref{sec:theoretical_foundations}, governs the physical setting of the work and dictates that the surrogate model must reconstruct three coupled state variables rather than a single field. The non-linear inter-porosity exchange term and the susceptibility of the system to viscous fingering make coarse-grid simulation alone insufficient and motivate the development of a learning-based reconstruction component.

Second, the analysis of numerical modeling in Section~\ref{sec:numerical_challenges} establishes the formal trade-off between truncation error and computational cost that justifies the surrogate approach: high-resolution simulation captures the relevant physics but is computationally prohibitive for ensemble workflows, while coarse-grid simulation is tractable but introduces a systematic loss of high-frequency content that classical interpolation cannot recover.

Third, the survey of evaluation metrics in Section~\ref{sec:metrics} demonstrates that no single metric is sufficient to characterize a hydrodynamic reconstruction. The evaluation framework adopted in the present work therefore combines a point-wise criterion (PSNR), a structural criterion (SSIM), two radiometric criteria adapted from remote sensing (ERGAS and SAM), a gradient-based criterion for front sharpness, and a physics-informed consistency criterion derived from the diffusive nature of pressure and from the geometric coherence of saturation fronts.

Fourth, the analysis of deep super-resolution architectures in Sections~\ref{sec:dl_architectures} and~\ref{sec:physics_sr} shows that the conditioned adversarial framework with physics-informed regularization is the most appropriate baseline for the target setting. Convolutional regression backbones, exemplified by MSRResNet, MDSR, and RFDN, establish a robust point-wise baseline but suffer from systematic smearing of fluid fronts. The adversarial framework recovers these fronts but introduces high-frequency artifacts unless complemented by physical regularization. Feature-wise linear modulation supplies the global operational context required to disambiguate plausible reconstructions. Alternative families based on recurrence, attention, or diffusion offer complementary advantages but fail to meet the interactive inference budget targeted in this work. The next chapter formalizes the resulting methodology and the associated experimental protocol.
